\chapter{Simulación}
\label{sec:simulacion}

\section{Circuito Simulado}

El circuito se simuló en KiCad con el motor ngspice (Anexo~\ref{sec:anexos}), utilizando los modelos
SPICE del LM335 y del MCP6001 incluidos en el anexo. El modelo del sensor posee un terminal $T$ cuya tensión
fija la temperatura en \si{\celsius}, por lo que se realizó un análisis DC barriendo la fuente $V_1$ de \SI{0}{\volt}
a \SI{50}{\volt} en pasos de \SI{1}{\volt} (\texttt{.dc V1 0 50 1}). El esquemático de simulación y la netlist completa se
encuentran en el anexo.

\section{Resultados de la Simulación}

La Figura~\ref{fig:sim_transferencia} muestra la tensión de salida en función de la temperatura junto a la recta
ideal, y la Tabla~\ref{tab:sim} resume los valores en puntos representativos. El error se expresa en grados
equivalentes, dividiendo la diferencia de tensión por la sensibilidad ideal de $\SI{100}{\milli\volt/\celsius}$.


\begin{table}[!ht]
  \centering
  \begin{tabular}{|c|c|c|c|c|}
    \hline
    $T$ [\unit{\celsius}] & $V_T$ [\unit{\volt}] & $V_o$ simulada [\unit{\volt}] & $V_o$ ideal [\unit{\volt}]
      & Error [\unit{\celsius}] \\
    \hline
    \num{0} & \num{2.7322} & \num{0.0486} & \num{0} & \num{0.49} \\
    \num{10} & \num{2.8321} & \num{1.0480} & \num{1} & \num{0.48} \\
    \num{20} & \num{2.9321} & \num{2.0475} & \num{2} & \num{0.47} \\
    \num{30} & \num{3.0320} & \num{3.0469} & \num{3} & \num{0.47} \\
    \num{40} & \num{3.1320} & \num{4.0463} & \num{4} & \num{0.46} \\
    \num{50} & \num{3.2319} & \num{4.9995} & \num{5} & \num{-0.01} \\
    \hline
  \end{tabular}
  \caption{Resultados de la simulación DC ($V_{REF} = \SI{2.7266}{\volt}$ simulada).}
  \label{tab:sim}
\end{table}

\begin{figure}[H]
  \centering
  \begin{tikzpicture}
  \colorlet{simazul}{blue!85!black}
  \colorlet{simrojo}{red!80!black}
  \begin{axis}[
    width=0.75\textwidth,
    height=4.6cm,
    grid=both,
    grid style={line width=.1pt, draw=gray!25},
    major grid style={line width=.2pt, draw=gray!45},
    xlabel={Temperatura del sensor [\unit{\celsius}]},
    ylabel={Tensión de salida $V_o$ [\unit{\volt}]},
    xmin=0, xmax=50,
    ymin=0, ymax=5.5,
    xtick={0,10,20,30,40,50},
    legend style={at={(0.5,1.03)}, anchor=south, font=\footnotesize, legend columns=2,
      /tikz/every even column/.append style={column sep=0.4cm}},
    ytick={0,1,2,3,4,5},
  ]
    \addplot[thick, color=simazul] table [x=temp, y=vout] {sources/sim/cas_dc.dat};
    \addlegendentry{$V_o$ simulada}
    \addplot[thick, dashed, color=simrojo, domain=0:50] {0.1*x};
    \addlegendentry{$V_o$ ideal}
  \end{axis}
\end{tikzpicture}

  \caption{Transferencia simulada del CAS en el rango de temperatura.}
  \label{fig:sim_transferencia}
\end{figure}

La sensibilidad obtenida entre \SI{0}{\celsius} y \SI{49}{\celsius} es de $\SI{99.94}{\milli\volt/\celsius}$,
prácticamente igual a la ideal, lo que confirma la ganancia $A_v = 10$ de la etapa restadora. La salida presenta un
desplazamiento aproximadamente constante de $\SI{48}{\milli\volt}$ (\SI{0.48}{\celsius}), originado en que el sensor
entrega $\SI{2.7322}{\volt}$ a \SI{0}{\celsius} (\SI{273.15}{\kelvin}) mientras que la referencia es de
$\SI{2.7266}{\volt}$: el divisor normalizado da $\SI{2.727}{\volt}$ en lugar de $\SI{2.73}{\volt}$, y a esto se suma la
tensión de \textit{offset} del seguidor. La diferencia de $\SI{5.5}{\milli\volt}$ queda amplificada diez veces.

Por encima de \SI{49.5}{\celsius} aproximadamente, la salida satura en $\SI{4.9995}{\volt}$ contra el riel positivo:
la recta simulada alcanzaría unos $\SI{5.05}{\volt}$ a \SI{50}{\celsius}. Por esto el error a \SI{50}{\celsius}
resulta casi nulo, pero se pierde la linealidad en el último medio grado del rango. Al mismo tiempo, la saturación
protege la entrada del conversor, ya que la salida nunca supera la alimentación.
